\chapter{Lung Transplant}

\section*{What Is This Surgery?}
Lung transplantation is a highly complex and definitive surgical procedure involving the replacement of one or both severely diseased lungs with healthy, functional lungs procured from a deceased donor. This life-saving intervention is meticulously reserved for patients suffering from end-stage pulmonary diseases that have become refractory to maximal medical and conventional surgical therapies, and which are associated with a high risk of mortality, typically within one to two years. The overarching goal of lung transplantation is to fundamentally restore adequate pulmonary function, thereby dramatically improving gas exchange efficiency, enhancing exercise capacity, and significantly elevating the patient's overall quality of life while extending their survival. It stands as a critical therapeutic option for conditions that otherwise carry a grim prognosis, but it necessitates a lifelong commitment to immunosuppressive therapy and rigorous, multidisciplinary medical follow-up.

\section*{Alternative Names}
\begin{itemize}
  \item Pulmonary Transplant
  \item Single Lung Transplant (SLT)
  \item Bilateral Lung Transplant (BLT)
  \item Double Lung Transplant
  \item Bilateral Sequential Lung Transplant
  \item Heart-Lung Transplant (for combined indications)
\end{itemize}

\section*{Relevant Surgical Anatomy}
A thorough understanding of thoracic anatomy is paramount for successful lung transplantation. The lungs themselves are paired organs situated within the pleural cavities, separated by the mediastinum. Each lung is divided into lobes: three on the right (upper, middle, lower) and two on the left (upper, lower, with a lingula). The pleura, a serous membrane, encases each lung, forming visceral and parietal layers. The hilum of each lung, located on its medial aspect, serves as the entry and exit point for the main bronchus, pulmonary artery, and pulmonary veins, collectively known as the lung root.

The tracheobronchial tree originates from the trachea, which bifurcates at the carina into the right and left main bronchi. These main bronchi further divide into lobar and segmental bronchi. The pulmonary circulation is unique: the pulmonary artery, originating from the right ventricle, carries deoxygenated blood to the lungs and divides into right and left main pulmonary arteries, which then branch extensively within the lung parenchyma. The pulmonary veins, typically two from each lung (superior and inferior), carry oxygenated blood back to the left atrium, often forming a common left atrial cuff for anastomosis. The bronchial arteries, usually arising from the aorta, supply the tracheobronchial tree and are typically not re-anastomosed during lung transplant, relying on collateral circulation from the pulmonary artery. Key neural structures include the vagus nerves, which descend alongside the trachea and esophagus, and the phrenic nerves, which run anterior to the lung hila, innervating the diaphragm. The recurrent laryngeal nerves, branches of the vagus, loop around the aorta (left) and subclavian artery (right) and ascend to the larynx; injury to these can cause vocal cord paralysis. Lymphatic drainage follows the bronchial and vascular structures to hilar and mediastinal lymph nodes.

\section{Surgical Landmarks:}
\begin{itemize}
  \item \textbf{Carina:} The bifurcation of the trachea into the main bronchi, a critical landmark for bronchial anastomosis.
  \item \textbf{Hila:} The central region of each lung where the main bronchus, pulmonary artery, and pulmonary veins enter/exit.
  \item \textbf{Pericardium:} The fibrous sac surrounding the heart, which must be carefully dissected to access the pulmonary veins and left atrium.
  \item \textbf{Aorta and Superior Vena Cava:} Major vessels in close proximity to the pulmonary arteries and veins, requiring careful dissection to avoid injury.
\end{itemize}

\section*{Surgical Principle \& Rationale}
The fundamental surgical principle of lung transplantation is the replacement of irreversibly damaged, non-functional lung parenchyma with healthy, viable donor lung tissue. End-stage lung diseases, regardless of etiology, culminate in progressive respiratory failure characterized by severe hypoxemia, hypercapnia, and often debilitating pulmonary hypertension, which can ultimately precipitate right ventricular failure. The transplanted lung or lungs are intended to provide immediate and sustained improvement in oxygenation and carbon dioxide elimination, thereby alleviating the profound symptoms of respiratory failure and improving systemic organ perfusion.

Technically, the procedure involves the meticulous creation of three primary anastomoses: the donor bronchus to the recipient's main bronchus, the donor pulmonary artery to the recipient's pulmonary artery, and the donor pulmonary veins (typically as a cuff of left atrium) to the recipient's left atrial cuff. The primary technical goals are to ensure patent, airtight airways, secure and hemostatic vascular connections for unimpeded blood flow, and to re-establish lymphatic and neural continuity to the extent surgically feasible, although the latter is often incomplete. While the transplanted lung immediately begins to function, it remains highly vulnerable to ischemia-reperfusion injury in the immediate post-operative period due to the period of cold preservation and subsequent re-establishment of blood flow. Furthermore, the lung is one of the most immunologically active organs, constantly exposed to environmental antigens, making it particularly susceptible to both acute and chronic rejection, necessitating aggressive and lifelong immunosuppressive therapy.

\section*{History \& Surgical Evolution}
The conceptualization of lung transplantation emerged in the mid-20th century, with pioneering experimental work in animal models laying the groundwork for human attempts. The first human lung transplant was performed by Dr. James Hardy and his team at the University of Mississippi Medical Center in 1963. The recipient, a patient with lung cancer, survived for only 18 days, succumbing to complications including renal failure and graft rejection, which underscored the immense immunological and technical challenges of the procedure at that time. Despite this initial limited success, the procedure unequivocally demonstrated the technical feasibility of lung replacement.

A pivotal turning point arrived with significant advancements in immunosuppressive therapy, most notably the introduction of cyclosporine in the early 1980s. This pharmacological breakthrough dramatically reduced the incidence and severity of acute rejection episodes. Leveraging these advancements, the Toronto Lung Transplant Group, under the leadership of Dr. Joel Cooper, achieved the first durable success with a single lung transplant in 1983. This was swiftly followed by their accomplishment of the first successful bilateral lung transplant in 1986. These landmark successes marked a critical inflection point, firmly establishing lung transplantation as a viable and effective treatment option for end-stage lung disease. Subsequent refinements in surgical techniques, including improved donor lung preservation, enhanced donor management protocols, and sophisticated post-operative care strategies, have continuously improved outcomes, transforming lung transplantation into a standard, albeit highly complex, therapeutic modality today.

\section*{Clinical Indications}
Lung transplantation is considered for patients with severe, irreversible lung disease who have a high risk of death within 1-2 years despite maximal medical therapy and who meet specific listing criteria.

\begin{itemize}
  \item End-stage Chronic Obstructive Pulmonary Disease (COPD) and emphysema, often exacerbated by alpha-1 antitrypsin deficiency, leading to severe airflow obstruction and hyperinflation.
  \item Idiopathic Pulmonary Fibrosis (IPF) and other progressive interstitial lung diseases, characterized by relentless scarring of the lung tissue and restrictive physiology.
  \item Cystic Fibrosis (CF) with severe bronchiectasis, recurrent infections, and progressive respiratory failure, often complicated by pulmonary hypertension.
  \item Primary Pulmonary Hypertension (PPH) or severe pulmonary hypertension secondary to other conditions, causing intractable right heart failure.
  \item Bronchiectasis not related to cystic fibrosis, leading to severe lung damage, recurrent infections, and respiratory compromise.
  \item Sarcoidosis with progressive pulmonary fibrosis and severe respiratory impairment refractory to medical management.
  \item Lymphangioleiomyomatosis (LAM), a rare progressive lung disease primarily affecting women, causing cystic lung destruction.
  \item Obliterative bronchiolitis or other forms of chronic rejection following a previous lung transplant, necessitating re-transplantation.
  \item Other rare conditions such as pulmonary veno-occlusive disease or Langerhans cell histiocytosis with end-stage lung involvement.
\end{itemize}

\section*{Clinical Positioning}
Lung transplantation is positioned as a primary, definitive, and often life-extending treatment for carefully selected patients with end-stage lung disease. It is emphatically not a first-line therapy but rather represents a last resort, indicated only after all other less invasive medical and conventional surgical management options have been thoroughly exhausted and the patient's prognosis remains demonstrably poor. Patients must undergo a rigorous, multidisciplinary evaluation to meet strict criteria for listing, which typically include a life expectancy of less than two years without transplant, the absence of any other viable medical or surgical alternatives, and a high likelihood of successful rehabilitation and unwavering adherence to a complex, lifelong post-transplant regimen. Therefore, while it serves as the ultimate and primary treatment for the underlying end-stage condition, its application is secondary to all other less invasive therapeutic modalities. In specific, acute scenarios, it may be considered a salvage procedure for patients experiencing rapid, life-threatening decompensation or catastrophic failure of a previous lung graft.

\section*{Surgical Technique \& Procedure}
Lung transplantation is a highly intricate surgical procedure performed under general anesthesia, often requiring advanced life support.

\begin{itemize}
  \item \textbf{Anesthesia and Cardiopulmonary Support:} The patient is placed under general anesthesia and intubated, typically with a double-lumen endotracheal tube to facilitate single-lung ventilation during the procedure. Comprehensive monitoring includes arterial and central venous lines, pulmonary artery catheterization, and transesophageal echocardiography (TEE). Cardiopulmonary bypass (CPB) or extracorporeal membrane oxygenation (ECMO) support is frequently utilized, especially during bilateral transplants or in patients with severe pulmonary hypertension or cardiac compromise, to maintain oxygenation and circulation while the diseased lungs are removed and the donor lungs implanted.
  \item \textbf{Patient Positioning and Incision:} For a single lung transplant (SLT), the patient is typically positioned in a lateral decubitus position, and a posterolateral thoracotomy incision is made on the side of the transplant. For bilateral lung transplantation (BLT), the patient is positioned supine, and a "clamshell" incision is often employed, which is a bilateral anterior thoracotomy with a transverse sternotomy, providing excellent exposure to both pleural cavities and the mediastinum.
  \item \textbf{Recipient Pneumonectomy:} The diseased recipient lung (or lungs) is meticulously dissected free from surrounding adhesions, which can be extensive in chronic lung disease. The pulmonary artery, pulmonary veins (often as a left atrial cuff), and main bronchus are identified, carefully clamped, divided, and prepared for the subsequent anastomoses.
  \item \textbf{Donor Lung Implantation:} The donor lung, which has been preserved in a cold solution (e.g., Perfadex), is then carefully brought into the recipient's chest cavity. The anastomoses are performed in a specific, critical sequence to minimize ischemia time and ensure optimal graft function:
    \begin{enumerate}
      \item \textbf{Bronchial Anastomosis:} The donor main bronchus is anastomosed to the recipient's main bronchus. This is a critical anastomosis, often reinforced with omentum or pericardial fat, as it is prone to complications like dehiscence, stenosis, or malacia.
      \item \textbf{Pulmonary Arterial Anastomosis:} The donor pulmonary artery is connected to the recipient's pulmonary artery, typically end-to-end.
      \item \textbf{Pulmonary Venous Anastomosis:} The donor pulmonary veins, usually as a cuff of left atrium, are anastomosed to the recipient's left atrial cuff.
    \end{enumerate}
  \item \textbf{Reperfusion and Airway Management:} Once all anastomoses are complete, blood flow is gradually restored to the donor lung. The lung is carefully re-inflated, and an air leak test is performed to ensure the integrity of the bronchial anastomosis. If CPB/ECMO was used, the patient is meticulously weaned off support.
  \item \textbf{Hemostasis and Closure:} Thorough hemostasis is achieved. Chest tubes are placed in the pleural space to drain fluid and air, facilitating lung re-expansion. The incision is then closed in layers, and the patient is transferred to the intensive care unit for close monitoring.
\end{itemize}

\section*{Preoperative Workup \& Preparation}
Preparation for lung transplantation is an extensive and multidisciplinary process, rigorously designed to ensure the patient is an optimal candidate and to minimize perioperative risks.

\begin{itemize}
  \item \textbf{Comprehensive Medical Evaluation:} This involves a rigorous assessment of all major organ systems, including detailed pulmonary function tests, cardiac stress testing, echocardiography, right heart catheterization, renal and hepatic function tests, and a thorough nutritional assessment.
  \item \textbf{Infection Screening:} Extensive screening for active or latent infections (e.g., viral hepatitis, HIV, CMV, EBV, fungal infections, tuberculosis, multi-drug resistant organisms) is crucial, as post-transplant immunosuppression can reactivate dormant pathogens or exacerbate existing ones.
  \item \textbf{Immunological Workup:} Human Leukocyte Antigen (HLA) typing and crossmatching with potential donors are performed to minimize the risk of hyperacute and acute rejection.
  \item \textbf{Psychosocial Assessment:} A thorough evaluation of the patient's psychological stability, social support system, and unwavering commitment to lifelong adherence to complex medical regimens is absolutely essential for long-term success.
  \item \textbf{Pulmonary Rehabilitation:} Many patients undergo intensive pre-transplant pulmonary rehabilitation to optimize their physical condition, respiratory muscle strength, and overall functional status, which can significantly improve post-operative recovery.
  \item \textbf{Bridging Therapies:} Patients may require advanced support such as supplemental oxygen, non-invasive ventilation, or even ECMO while awaiting a suitable donor lung.
  \item \textbf{Pre-operative Instructions:} Patients are typically instructed to be NPO (nothing by mouth) for a specified period before surgery. Essential medications may be adjusted or held, and prophylactic antibiotics are administered intravenously prior to incision.
  \item \textbf{Informed Consent:} Detailed discussions regarding the significant risks, potential benefits, available alternatives, and lifelong implications of transplantation are conducted with the patient and family, and comprehensive informed consent is obtained.
\end{itemize}

\section*{Contraindications \& Precautions}
\textbf{Absolute Contraindications:}
\begin{itemize}
  \item Active malignancy within the last 5 years (excluding treated non-melanoma skin cancers or carcinoma in situ).
  \item Untreatable active infection (e.g., HIV with uncontrolled viral load, active tuberculosis, multi-drug resistant bacterial or fungal infections that cannot be eradicated).
  \item Severe, irreversible dysfunction of other major organ systems (e.g., severe renal failure not amenable to dialysis/transplant, severe liver cirrhosis, untreatable coronary artery disease, or severe cerebrovascular disease).
  \item Significant uncorrectable chest wall or spinal deformity that would severely compromise lung expansion or surgical access.
  \item Documented non-adherence to medical therapy, active substance abuse (alcohol, illicit drugs), or active smoking within a specified period (typically 6 months).
  \item Severe, irreversible neurological impairment that would preclude rehabilitation or adherence to post-transplant care.
  \item Uncontrolled psychiatric illness that would interfere with compliance or recovery.
  \item Profound cachexia or morbid obesity (BMI >40 kg/m\textasciicircum2) that significantly increases surgical risk.
\end{itemize}

\textbf{Relative Contraindications and Precautions:}
\begin{itemize}
  \item Age greater than 65-70 years, depending on overall physiological status and comorbidities.
  \item Moderate obesity (BMI 30-35 kg/m\textasciicircum2).
  \item Critical colonization with highly resistant organisms (e.g., Burkholderia cepacia complex in cystic fibrosis), which increases infection risk.
  \item Significant osteoporosis with vertebral compression fractures or other severe bone disease.
  \item History of recurrent aspiration or severe, uncontrolled gastroesophageal reflux disease (GERD), which can contribute to chronic rejection.
  \item Prior extensive thoracic surgery or pleurodesis, which can significantly increase surgical complexity and operative time due to adhesions.
  \item Lack of adequate psychosocial support system or financial resources to manage lifelong post-transplant care.
  \item Untreated or poorly controlled diabetes mellitus.
\end{itemize}

\section*{Complications \& Risks}
Lung transplantation carries significant risks, both immediate and long-term, due to the complexity of the surgery, the fragility of the graft, and the necessity for lifelong immunosuppression.

\begin{itemize}
  \item \textbf{General Surgical Complications:}
    \begin{itemize}
      \item \textbf{Bleeding:} Intraoperative hemorrhage or significant post-operative bleeding requiring transfusion or surgical re-exploration.
      \item \textbf{Infection:} Surgical site infection, nosocomial pneumonia, urinary tract infection, or systemic sepsis, particularly in an immunocompromised patient.
      \item \textbf{Anesthesia Complications:} Adverse reactions to anesthetic agents, cardiac events (e.g., myocardial infarction, arrhythmias), stroke, or prolonged respiratory depression.
      \item \textbf{Thromboembolic Events:} Deep vein thrombosis (DVT) and pulmonary embolism (PE).
    \end{itemize}
  \item \textbf{Procedure-Specific Risks:}
    \begin{itemize}
      \item \textbf{Primary Graft Dysfunction (PGD):} A severe form of acute lung injury occurring within the first 72 hours post-transplant, characterized by pulmonary edema and impaired gas exchange, ranging from mild to severe, and is a leading cause of early mortality.
      \item \textbf{Acute Rejection:} Occurs in the early post-transplant period, mediated by the recipient's immune system attacking the donor lung, requiring intensified immunosuppression.
      \item \textbf{Chronic Rejection (Bronchiolitis Obliterans Syndrome - BOS):} The leading cause of late morbidity and mortality, characterized by progressive, irreversible airflow obstruction due to inflammation and fibrosis of the small airways.
      \item \textbf{Anastomotic Complications:}
        \begin{itemize}
          \item \textbf{Bronchial Anastomosis:} Dehiscence (separation), stenosis (narrowing), malacia (softening), or fistula formation of the airway, potentially leading to infection, airway collapse, or pneumothorax.
          \item \textbf{Vascular Anastomosis:} Thrombosis (clot formation) or stenosis of the pulmonary artery or veins, compromising blood flow to/from the graft, potentially leading to graft failure.
        \end{itemize}
      \item \textbf{Immunosuppression-Related Complications:}
        \begin{itemize}
          \item Increased susceptibility to opportunistic infections (viral, fungal, bacterial, protozoal).
          \item Renal dysfunction or failure due to calcineurin inhibitors.
          \item Hypertension and new-onset diabetes mellitus.
          \item Post-transplant lymphoproliferative disorder (PTLD) and other malignancies (e.g., skin cancers, Kaposi's sarcoma).
          \item Osteoporosis and bone fractures.
          \item Gastrointestinal side effects (nausea, diarrhea).
        \end{itemize}
      \item \textbf{Phrenic Nerve Injury:} Can lead to diaphragmatic paralysis and significantly impaired respiratory mechanics, especially if bilateral.
      \item \textbf{Recurrent Laryngeal Nerve Injury:} Can cause vocal cord paralysis, leading to dysphonia and increased risk of aspiration.
      \item \textbf{Gastroesophageal Reflux Disease (GERD):} Often exacerbated post-transplant, contributing to chronic rejection and aspiration risk.
      \item \textbf{Pulmonary Edema:} Non-cardiogenic pulmonary edema can occur due to reperfusion injury or lymphatic disruption.
    \end{itemize}
\end{itemize}

\section*{Postoperative Recovery Timeline}
The postoperative period following lung transplantation is exceptionally critical and demands intensive, specialized care within a dedicated transplant intensive care unit (ICU). Immediately after surgery, the patient remains intubated and on mechanical ventilation, often with advanced hemodynamic support. Close monitoring of vital signs, hemodynamics, oxygenation, chest tube output, and neurological status is paramount. Aggressive pain management is initiated, and immunosuppressive medications are promptly started to prevent acute rejection. Early mobilization, often with the assistance of physical therapists and respiratory therapists, is strongly encouraged as soon as the patient is hemodynamically stable to prevent deconditioning, atelectasis, and other pulmonary complications.

Over the first 24-48 hours, the primary focus is on gradually weaning the patient from mechanical ventilation, managing chest tubes (which are typically removed when drainage is minimal and no air leak is present), and implementing aggressive pulmonary hygiene strategies to prevent atelectasis and pneumonia. Surveillance bronchoscopies may be performed within the first few days to assess the integrity of the bronchial anastomosis and obtain samples for infection surveillance. As the patient stabilizes and is successfully extubated, they are transferred out of the ICU, typically within the first week. During the first 1-2 weeks, the emphasis shifts to intensive physical therapy, comprehensive patient education regarding their complex medication regimen, and vigilant monitoring for early complications such as infection or acute rejection. Diet is gradually advanced from clear liquids to a regular diet as tolerated. Long-term recovery involves a commitment to lifelong immunosuppression, regular follow-up appointments, diagnostic tests, and adherence to a healthy lifestyle to maximize graft survival and overall well-being, with full functional recovery often taking several months.

\section*{Surgical Findings \& Pathology}
During lung transplantation, the surgeon meticulously documents the intraoperative findings, which are crucial for understanding the patient's disease burden and guiding postoperative management. These findings typically include the extent of adhesions (e.g., dense pleural adhesions from prior infections or pleurodesis), the degree of emphysema, fibrosis, or bronchiectasis in the native lung, the presence and severity of pulmonary hypertension, and the condition of the mediastinal structures and heart. Any unexpected findings, such as occult malignancy or significant anatomical variations, are also noted. The quality and size match of the donor lung are also assessed visually and by palpation.

The resected native lung tissue is sent for comprehensive pathological examination. The pathology report confirms the underlying diagnosis (e.g., severe emphysema, idiopathic pulmonary fibrosis, cystic fibrosis with extensive bronchiectasis and abscess formation, pulmonary hypertension with vascular remodeling). It details the extent of disease, the presence of any secondary complications such as infection, malignancy, or amyloidosis, and provides a definitive histological assessment. For the donor lung, a pre-implantation biopsy may be performed to assess its health and rule out occult disease. Post-transplant, surveillance transbronchial biopsies are routinely performed to monitor for acute cellular rejection, antibody-mediated rejection, and signs of chronic rejection (bronchiolitis obliterans syndrome), as well as to detect opportunistic infections within the graft. The correlation of surgical findings with the pathological report provides a complete picture of the disease state and the success of its removal.

\section*{Expected Postoperative Course}
A normal, successful postoperative course following lung transplantation is characterized by a progressive and steady improvement in respiratory function and overall well-being. Patients are expected to be weaned from mechanical ventilation within the first 24-72 hours, though this can vary based on individual factors and the complexity of the surgery. Pain, initially managed aggressively with intravenous analgesics, gradually subsides and transitions to oral medications. Chest tubes are typically removed within 3-7 days as lung re-expansion is complete and drainage diminishes. Patients are encouraged to ambulate early and participate actively in physical therapy, which is crucial for regaining strength and lung capacity. The diet is advanced as tolerated, usually starting with clear liquids and progressing to a regular diet. The primary goal is to achieve stable graft function, evidenced by improving spirometry values and resolution of pre-transplant symptoms like dyspnea. Patients are typically discharged from the hospital within 2-4 weeks, provided they are medically stable, have adequate pain control, are able to ambulate, and demonstrate proficiency in managing their complex medication regimen.

\section*{Postoperative Care \& Follow-up}
Lifelong, rigorous postoperative care and follow-up are absolutely essential after lung transplantation to monitor for complications, manage immunosuppression, and ensure optimal graft function and patient well-being.

\begin{itemize}
  \item \textbf{Regular Clinic Visits:} Frequent visits to the transplant center are scheduled, initially weekly, then monthly, and eventually annually, for comprehensive clinical assessment, medication management, and psychosocial support.
  \item \textbf{Pulmonary Function Tests (PFTs):} Regular spirometry, often performed daily at home and during clinic visits, is crucial to detect subtle changes in lung function that may signal rejection or infection.
  \item \textbf{Bronchoscopy with Transbronchial Biopsy and BAL:} Periodic procedures (e.g., at 1, 3, 6, 12 months post-transplant, and then as clinically indicated) are performed to screen for acute and chronic rejection and to identify infections.
  \item \textbf{Imaging Studies:} Routine chest X-rays and occasional high-resolution CT scans of the chest are utilized to monitor for infection, rejection, effusions, or other pulmonary complications.
  \item \textbf{Blood Tests:} Regular monitoring of immunosuppressant drug levels, renal and hepatic function, complete blood counts, and infection markers (e.g., CMV viral load) is critical.
  \item \textbf{Physical Therapy and Pulmonary Rehabilitation:} Ongoing participation in rehabilitation programs is vital to maintain physical conditioning, optimize lung function, and improve overall quality of life.
  \item \textbf{Vaccinations:} Adherence to a specific vaccination schedule (e.g., influenza, pneumococcal, COVID-19) is crucial to protect against preventable infections in an immunocompromised state.
  \item \textbf{Cancer Screening:} Increased surveillance for post-transplant malignancies, especially skin cancers, PTLD, and other solid organ tumors, is recommended.
  \item \textbf{Medication Adherence:} Strict adherence to the lifelong immunosuppressive regimen and other prescribed medications is paramount to prevent rejection and manage comorbidities.
\end{itemize}
Patients are thoroughly educated on warning signs such as fever, increasing cough, shortness of breath, decreased exercise tolerance, or general malaise, and instructed to immediately contact their transplant team if these symptoms arise, as they may indicate infection or rejection.

\section*{Epidemiology}
Lung transplantation remains a relatively rare but increasingly common procedure for end-stage lung disease globally. Approximately 4,000 to 5,000 lung transplants are performed annually worldwide, with the United States and Europe accounting for the majority of these procedures. The typical recipient age ranges from 50 to 65 years, though pediatric and older adult transplants are also performed under specific, carefully evaluated circumstances. There is a slight male predominance among recipients. The most common indications for lung transplantation include Chronic Obstructive Pulmonary Disease (COPD) and emphysema (approximately 25-30\% of cases), Idiopathic Pulmonary Fibrosis (IPF) (20-25\%), and Cystic Fibrosis (CF) (15-20\%). Other indications like alpha-1 antitrypsin deficiency, pulmonary hypertension, and sarcoidosis account for smaller percentages. Despite significant advancements in surgical technique and immunosuppression, lung transplantation carries substantial morbidity and mortality; 30-day mortality rates typically range from 5-10\%, primarily due to primary graft dysfunction, infection, or surgical complications. The demand for donor lungs consistently outstrips supply, leading to significant waiting lists and patient mortality while awaiting transplant.

\section*{Outcomes \& Prognosis}
The outcomes of lung transplantation have steadily improved over the past decades, offering significant survival and quality of life benefits for carefully selected patients. According to data from the International Society for Heart and Lung Transplantation (ISHLT) Registry, the median survival after lung transplantation is approximately 6-7 years for bilateral lung transplants and 4-5 years for single lung transplants. One-year survival rates typically range from 85-90\%, with 5-year survival rates around 50-60\%. Functional outcomes are generally excellent, with the vast majority of recipients experiencing a dramatic improvement in exercise capacity, a significant reduction in dyspnea, and a substantial enhancement in their overall quality of life, often returning to activities of daily living and even light work. However, chronic rejection, primarily manifesting as Bronchiolitis Obliterans Syndrome (BOS), remains the leading cause of long-term morbidity and mortality, affecting about 50\% of patients by 5 years post-transplant. Other factors influencing prognosis include the underlying disease (e.g., CF patients often have better long-term survival than IPF patients), the severity of primary graft dysfunction, the recipient's age, and critically, adherence to the lifelong immunosuppressive regimen. Despite these persistent challenges, lung transplantation offers a viable and often life-extending option for those facing otherwise terminal lung disease, as evidenced by ongoing improvements in registry data.

\section*{References}
\begin{enumerate}
  \item MedlinePlus. Lung transplant. U.S. National Library of Medicine; 2024. Available from: https://medlineplus.gov/lungtransplant.html
  \item Kotloff RM, et al. Official American Thoracic Society/European Respiratory Society/International Society for Heart and Lung Transplantation Guidelines for the Recipient Workup, Donor Selection, and Medical and Surgical Management of Lung Transplantation. Am J Respir Crit Care Med. 2021;204(1):e1-e64.
  \item Patterson GA, et al. Lung Transplantation. In: Shields TW, et al., editors. Shields' General Thoracic Surgery. 8th ed. Wolters Kluwer; 2018. p. 1157-1184.
  \item International Society for Heart and Lung Transplantation (ISHLT). ISHLT Registry for Heart and Lung Transplantation. Available from: https://ishlt.org/registries
  \item Sabiston DC, et al. Sabiston Textbook of Surgery: The Biological Basis of Modern Surgical Practice. 21st ed. Elsevier; 2021.
\end{enumerate}

\clearpage